% Lecture 07: Наследование и полиморфизм

\section{Лекция 7. Наследование и полиморфизм}
\label{sec:lecture-07}

В блоках~5--6 один class описывал один самостоятельный тип: состояние,
инвариант, создание, копирование и уничтожение. Теперь несколько видов
проверок блока курса должны запускаться одинаково, хотя каждая проверяет своё
условие. Нужен общий контракт поведения, но не обязательно одна реализация.

\begin{importantbox}[Главный вопрос]
Когда один тип --- разновидность другого, как C++
выбирает реализацию операции runtime и почему копирование derived object в
base object меняет смысл программы?
\end{importantbox}

\subsection{Inheritance --- отношение между типами}

\begin{cppcode}
class Report {
public:
    virtual std::string kind() const { return "Report"; }
};

class LectureReport : public Report {
public:
    std::string kind() const override { return "LectureReport"; }
};
\end{cppcode}

\texttt{Report} --- base class, \texttt{LectureReport} --- derived class.
Запись \texttt{: public Report} сообщает design-намерение: объект
\texttt{LectureReport} допустимо использовать там, где контракт требует
\texttt{Report}.

\begin{definition}
\textbf{Inheritance} вводит отношение base--derived между типами.
\textbf{Runtime polymorphism} позволяет обращаться к объектам разных derived
types через общий base interface и выбирать virtual implementation по типу
фактического объекта.
\end{definition}

Повтор кода сам по себе недостаточен. Если два несвязанных типа имеют функцию
\texttt{print}, один не становится разновидностью другого. Общий алгоритм
можно вынести в функцию или member-object. Inheritance оправдано, когда важны
общий контракт и substitutability.

\subsection{Is-a и has-a до синтаксиса}

Сначала формулируем связь предметной области:
\texttt{Car is-a Vehicle}: автомобиль --- вид транспортного средства; public inheritance может соответствовать смыслу; \texttt{Car has-a Engine}: автомобиль имеет двигатель; engine --- member/collaborator, то есть нужна composition; \texttt{CourseAudit has checking components}: аудит использует компоненты проверок, но сам не совпадает с конкретной проверкой.

Интуитивная проверка подстановки: если функция обещает работать с любым
\texttt{Vehicle}, сохраняет ли \texttt{Car} ожидания base contract? Это
подводит к Liskov substitution, но полный SOLID относится к блоку~8.

Неверный пример --- \texttt{class Engine : public Car}: двигатель --- не автомобиль. Наследование ради доступа к готовому коду создало бы ложную
связь. Аналогично \texttt{CourseAudit : BlockCheck} сомнительно, если аудит
координирует несколько checks, а не представляет одну взаимозаменяемую check.

\begin{whybox}[Composition не «побеждает всегда»]
Для has-a/uses-a и переиспользования реализации без substitutability обычно
предпочтительна composition. Для independently extensible операций, которые
должны вызываться через единый interface runtime, public inheritance может
точнее выразить задачу.
\end{whybox}

\subsection{Base subobject и доступ к members}

Derived object содержит base class subobject. Это корректная модель языка:
при создании \texttt{LectureReport} сначала инициализируется его
\texttt{Report} subobject, затем собственные members derived. Не следует
говорить «base просто лежит как поле»: у base subobject специальные правила,
включая conversions и virtual dispatch.

\begin{cppcode}
class Report {
public:
    const std::string& name() const { return name_; }
protected:
    void mark_checked() { checked_ = true; }
private:
    std::string name_;
    bool checked_{};
};
\end{cppcode}

Derived может вызвать \texttt{mark\_checked}, но не обратиться прямо к
\texttt{name\_} или \texttt{checked\_}. Private state существует в base
subobject, но base сохраняет контроль над invariant. Много protected data
ослабляет encapsulation: каждый derived сможет создавать новые сочетания
state. Обычно private data и узкие protected/public operations проще
контролировать.

\subsection{Public, protected и private inheritance}

\begin{center}\small
\begin{tabularx}{\textwidth}{@{}lXXX@{}}
\toprule
Форма & public member base & protected member base & обычный смысл \\
\midrule
\texttt{public Base} & public & protected & is-a, substitutability \\
\texttt{protected Base} & protected & protected & специальная реализация \\
\texttt{private Base} & private & private & implementation technique \\
\bottomrule
\end{tabularx}
\end{center}

Private members base всё равно не становятся доступны derived напрямую.
Protected/private inheritance существуют, но не выражают публичное is-a для
caller. Основной production-style лекции --- \texttt{public Base}. Access
inheritance не надо путать с access sections внутри class.

\subsection{Static type и dynamic type --- центр модели}

\begin{cppcode}
LectureReport lecture;
Report& reference{lecture};
Report* pointer{&lecture};
\end{cppcode}

\begin{definition}
\textbf{Static type} expression известен compiler из исходного кода. Static
types имён \texttt{reference} и \texttt{pointer} --- \texttt{Report\&} и
\texttt{Report*}. \textbf{Dynamic type} обозначенного ими most-derived object
во время вызова --- \texttt{LectureReport}.
\end{definition}

Reference и pointer здесь не владеют объектом. \texttt{lecture} имеет
automatic storage duration и должен оставаться жив дольше их использования.
Pointer может быть \texttt{nullptr}; reference обозначает обязательный object.

\begin{cppcode}
Report report = lecture;
\end{cppcode}

У самостоятельного \texttt{report} static и dynamic object type ---
\texttt{Report}. Это новый base object, не скрытый derived object.

\subsection{Non-virtual и virtual dispatch}

\begin{cppcode}
class Report {
public:
    std::string non_virtual_kind() const { return "Report"; }
    virtual std::string virtual_kind() const { return "Report"; }
};

Report& view{lecture};
view.non_virtual_kind(); // Report: static dispatch
view.virtual_kind();     // LectureReport: dynamic dispatch
\end{cppcode}

Non-virtual member выбирается по static type expression. Для virtual call
через base pointer/reference final overrider выбирается по dynamic type
object runtime. Virtual member остаётся virtual в derived без повторного
слова \texttt{virtual}. Modern C++ всё равно пишет:

\begin{cppcode}
std::string virtual_kind() const override;
\end{cppcode}

\texttt{override} не «включает virtual заново». Оно фиксирует намерение и
требует от compiler найти подходящую virtual function base. Если забыть
\texttt{const}, изменить parameter type или ошибиться в имени, compilation
завершится ошибкой вместо создания другой функции.

\begin{mistakebox}[Overriding не overloading]
Overriding заменяет virtual behavior base для подходящей signature в derived.
Overloading создаёт функции одного имени с разными parameter lists. Derived
declaration с похожим именем способно скрыть base overloads; совпадения имени
недостаточно. \texttt{override} защищает от этой путаницы.
\end{mistakebox}

\subsection{Object slicing: новое base value}

\begin{cppcode}
void print_by_value(Report report);
LectureReport lecture{"lecture 07", 500};
print_by_value(lecture);
\end{cppcode}

Parameter \texttt{report} --- новый \texttt{Report} object. Он
инициализируется из base subobject argument; derived-specific members вроде
\texttt{lines\_} в него не входят. Dynamic type parameter равен
\texttt{Report}, поэтому даже virtual call выбирает base behavior.

\begin{definition}
\textbf{Object slicing} --- потеря derived-specific части при создании base
object by value из derived object. Копируется base portion/state, но новый
object --- не alias исходного derived object.
\end{definition}

\begin{cppcode}
void print_by_reference(const Report& report);
print_by_reference(lecture); // нет нового Report, slicing отсутствует
Report* view{&lecture};       // non-owning view того же object
\end{cppcode}

Reference/pointer обозначает существующий \texttt{LectureReport}, поэтому
dynamic type сохраняется. Polymorphic objects обычно передают по reference или
pointer, если нужна identity/dynamic behavior, а не новая base value. Slicing
не обязан вызывать warning: программа может быть формально корректна, но иметь
неверную semantics.

\subsection{Pure virtual function и abstract class}

\begin{cppcode}
class BlockCheck {
public:
    virtual ~BlockCheck() = default;
    virtual bool passes(const CourseBlock& block) const = 0;
};
\end{cppcode}

\texttt{= 0} делает функцию pure virtual. Class с pure virtual function
--- abstract: создать самостоятельный \texttt{BlockCheck} нельзя.
Concrete derived предоставляет final overrider:

\begin{cppcode}
class LectureExistsCheck final : public BlockCheck {
public:
    bool passes(const CourseBlock& block) const override {
        return block.lecture_exists();
    }
};
\end{cppcode}

Abstract class задаёт interface-like contract, но C++ не требует превращать
все classes в «интерфейсы». Он способен иметь state, implemented members и
constructors. Здесь stateless interface точнее задачи. \texttt{final}
запрещает дальнейшее наследование от concrete check; это локальное решение.

\subsection{Virtual destructor и polymorphic deletion}

Главная lifetime-опасность возникает не от любого base pointer, а от
уничтожения объекта через него:

\begin{cppcode}
class SafeBase {
public:
    virtual ~SafeBase() = default;
};
class SafeDerived : public SafeBase {
public:
    ~SafeDerived() override = default;
};

SafeBase* base = new SafeDerived;
delete base; // destruction начинается с SafeDerived
\end{cppcode}

Virtual destructor выбирается по dynamic type. Сначала выполняется destructor
derived class, затем уничтожаются derived members и base subobject через
destructor base class. Если base используется
polymorphically и контракт допускает \texttt{delete} через base pointer,
destructor base должен быть virtual.

\begin{ubbox}[Non-virtual destructor]
Если static type operand \texttt{delete} отличается от dynamic type object,
для корректности base должен иметь virtual destructor с учётом точных правил
языка. Удаление \texttt{UnsafeDerived} через \texttt{UnsafeBase*} без него
имеет undefined behavior. Это не гарантирует ни «только base destructor», ни
crash: после UB стандарт не задаёт результат.
\end{ubbox}

Raw \texttt{new}/\texttt{delete} используются только в изолированном negative
experiment. Основной design применяет automatic objects и non-owning views с
очевидным lifetime. Smart pointers появятся позже.

\subsection{Construction и destruction иерархии}

\begin{cppcode}
LectureReport::LectureReport(const std::string& name, int lines)
    : Report{name}, lines_{lines} {}
\end{cppcode}

Base subobject создаётся до members derived, а тело derived constructor
выполняется после initialization. Уничтожение идёт обратно: сначала derived
destructor и members derived, затем base. Это продолжает lifetime model
блока~6.

Virtual calls из constructors/destructors требуют осторожности. Во время base
construction более-derived часть ещё не построена, а во время base
destruction уже уничтожена. Dispatch ограничен текущим этапом иерархии, а не
ведёт себя как вызов через полностью готовый most-derived object. Не следует
строить initialization/cleanup на ожидании derived override из base
constructor/destructor.

\subsection{Composition и inheritance на одной задаче}

Проект проверяет \texttt{CourseBlock}: есть ли lecture file, examples
directory и содержимое. В inheritance-design каждая проверка is-a
\texttt{BlockCheck}:

\begin{cppcode}
LectureExistsCheck lecture;
ExamplesExistCheck examples;
ContentCheck content;
const std::array<const BlockCheck*, 3> checks{
    &lecture, &examples, &content
};
for (const BlockCheck* check : checks) {
    results.push_back(check->run(block));
}
\end{cppcode}

Checks живут до конца функции, поэтому pointers --- безопасные non-owning
views. Плюс: caller одинаково запускает independently extensible types, не
зная concrete class. Цена: base interface, virtual functions и class на
каждое правило создают coupling и boilerplate.

Composition-design использует collaborators:

\begin{cppcode}
class CourseAudit {
public:
    std::vector<CheckResult> run(const CourseBlock& block) const;
private:
    PathChecks paths_;
    ContentChecks content_;
};
\end{cppcode}

\texttt{CourseAudit has-a PathChecks} и \texttt{ContentChecks}. Для
фиксированного локального набора правил форма короче, вызовы static и нет
ложной is-a relation. Если проверки должны независимо выбираться runtime через
стабильный contract, polymorphic \texttt{BlockCheck} полезнее. Решение зависит
от отношения и точки изменения, а не от лозунга.

\subsection{Memory model и стоимость без мифов}

Стандарт задаёт observable behavior virtual calls, но не требует vtable или
vptr layout. Обычная ABI implementation хранит в polymorphic object
служебный pointer (vptr) на таблицу virtual functions (vtable), а call может
стать indirect call. Это mental model, не переносимая layout guarantee.

\texttt{sizeof(Report)} и \texttt{sizeof(LectureReport)} можно измерить, но
числа зависят от compiler, ABI, members, alignment и platform. Metadata в
обычной реализации способна увеличить object size. Inheritance без virtual
functions само по себе не обязано добавлять runtime dispatch overhead.

\begin{performancebox}[Virtual не означает «всегда медленно»]
Virtual call часто реализуется как indirect call, что влияет на inlining и
branch prediction. Compiler иногда знает dynamic type и devirtualizes вызов.
Итог зависит от optimization/context; allocations, algorithm и locality могут
быть важнее одного dispatch. Нужны benchmark/assembly конкретной сборки.
\end{performancebox}

Dynamic dispatch --- выбранная стоимость runtime substitutability. Когда
вариативность не нужна, ordinary function или composition оставляют выбор
compile time. Цель --- платить за механизм там, где contract требует runtime
choice.

\subsection{Эксперимент 1: polymorphism\_failure\_lab}

Safe executable показывает контраст:

\begin{outputcode}
static dispatch: Report
dynamic dispatch: LectureReport
pointer dispatch: LectureReport
sliced value: kind=Report, name=07_inheritance_and_polymorphism
reference: kind=LectureReport, name=07_inheritance_and_polymorphism
virtual destructor trace:
destroy SafeDerived
destroy SafeBase
POLYMORPHISM_FAILURE_LAB=PASS
\end{outputcode}

Полный output также подтверждает protected operation, derived state и concrete
реализацию abstract interface. Три CTest cases требуют compilation failure:
нет \texttt{const} у \texttt{override}; создаётся abstract base; derived
напрямую обращается к private base state.

\begin{shellcode}
cd examples/07_inheritance_and_polymorphism/01_polymorphism_failure_lab
cmake -S . -B build -DCMAKE_BUILD_TYPE=Debug
cmake --build build
./build/polymorphism_failure_lab
ctest --test-dir build --output-on-failure
\end{shellcode}

Unsafe executable не входит в normal run:

\begin{shellcode}
cmake -S . -B build-sanitize -DENABLE_UNSAFE_CASE=ON \
  -DCMAKE_BUILD_TYPE=Debug
cmake --build build-sanitize
ctest --test-dir build-sanitize --output-on-failure
\end{shellcode}

На фактической GCC~13/ASan сборке unsafe process завершился nonzero, а test
нашёл relevant AddressSanitizer diagnostic. PASS означает выполненный bad path
и ожидаемую причину, не «UB стало безопасным». На другой toolchain diagnostic
может отличаться; языковое правило остаётся.

\subsection{Эксперимент 2: project design lab}

\texttt{02\_inheritance\_vs\_composition\_project\_lab} читает реальные
lecture/examples блока~07. Оба design дают одинаковый результат:

\begin{outputcode}
inheritance:
  lecture exists=PASS
  examples directory exists=PASS
  lecture and examples non-empty=PASS
composition:
  lecture exists=PASS
  examples directory exists=PASS
  lecture and examples non-empty=PASS
same observable results=true
PROJECT_DESIGN_LAB=PASS
\end{outputcode}

\begin{shellcode}
cd examples/07_inheritance_and_polymorphism/02_inheritance_vs_composition_project_lab
cmake -S . -B build -DCMAKE_BUILD_TYPE=Debug
cmake --build build
./build/project_lab --root ../../..
ctest --test-dir build --output-on-failure
\end{shellcode}

PASS требует exit code~0, трёх успешных checks и эквивалентного observable
result. Для fixed audit composition проще. Inheritance-version показывает
полезную точку: новый concrete \texttt{BlockCheck} можно передать тому же
caller. Это reusable audit idea, но не преждевременный framework.

\subsection{Сильные стороны и цена механизма}

C++ включает runtime polymorphism в type system: compiler проверяет
conversions и \texttt{override}, programmer выбирает static/dynamic dispatch,
а public inheritance выражает настоящую substitutability. Composition остаётся
отдельным инструментом, и virtual abstraction не навязывается каждому class.

Цена ошибок: ложная hierarchy усиливает coupling; slicing может тихо изменить
semantics; non-virtual destructor опасен при polymorphic deletion; layout
становится менее очевидным; indirect dispatch может влиять на optimization.
Inheritance ради одного code reuse часто ухудшает design, потому что обещает
отношение между типами, которого в задаче нет.

\subsection{Проверка понимания}

После блока нужно уметь:
1) назвать base/derived и прочитать \texttt{class D : public B}; 2) отличить is-a от has-a и предложить composition для второго; 3) объяснить public inheritance и осторожное применение protected; 4) определить static/dynamic type до объяснения virtual call; 5) предсказать non-virtual и virtual output через base view; 6) объяснить \texttt{override} как compiler-checked intent; 7) причинно описать slicing и отличие base value от alias; 8) объяснить pure virtual function и abstract class; 9) обосновать virtual destructor именно polymorphic deletion; 10) назвать порядок base/derived construction/destruction; 11) сравнить inheritance/composition без абсолютного правила; 12) отделить vptr/vtable model от standard guarantee; 13) оценивать virtual cost по контексту, а не словом «всегда».

\subsection{Источники для сверки}

Base/derived classes, virtual functions, final overrider, abstract classes,
initialization/destruction и delete-expression сверены с C++17 working draft
N4659, прежде всего clauses~10, 12 и 15. Педагогическая последовательность
сопоставлена с MIT 6.096 Lecture~7, а classes/lifetime --- с Lectures~6 и~8.
Современное применение \texttt{override} сверено с локальным \textit{Modern
C++ Tutorial}. Vtable/vptr названы обычной ABI strategy, не требованием
стандарта; performance context согласован с локальным C++ Performance TR.
